\clearpage\thispagestyle{empty}
\pdfbookmark[0]{Quellen zu Komet (April 1813)}{comet1813apr-sources}
\begin{center}\textbf{Quellen zu Komet (April 1813)}\end{center}
\noindent Diese KI-Übertragung ins heutige Deutsch folgt dem Auszug mit der Überschrift Comet in den \href{https://gdz.sub.uni-goettingen.de/id/PPN236007467}{\textit{Werken}, Band~VI} von Gauß (Göttingen, 1874), S.~364--365. Er stammt aus Nr.~XLIV der Zeitschrift \textit{Monatliche Correspondenz zur Beförderung der Erd- und Himmelskunde}, Band~XXVII (April 1813), S.~\href{https://api.digitale-sammlungen.de/iiif/image/v2/bsb10538616_00402/full/1600,/0/default.jpg}{388}, \href{https://api.digitale-sammlungen.de/iiif/image/v2/bsb10538616_00403/full/1600,/0/default.jpg}{389}. Die Links führen zum digitalisierten Exemplar der Bayerischen Staatsbibliothek. Der Auszug beginnt mit den fünf von Gauß bestimmten Positionen und endet damit, dass Enke an der Verbesserung der Elemente arbeitet.
\par\smallskip\noindent Sowohl in der Werkausgabe als auch im Erstdruck stehen in Hardings Beobachtung 71 Sekunden. In Gauß’ zweitem Satz von Elementen gibt die Werkausgabe für den Periheldurchgang den 15.93437 Mai an, der Erstdruck den 19,93437 Mai. Die vier Kometenzeichen folgen der Werkausgabe. KI-Übertragung und KI-Satz.
\clearpage
